\chapter{Validación y discusión}

\section{Validación del modelo numérico}

La consistencia del modelo computacional se evalúa mediante la comparación de sus resultados con las expresiones analíticas utilizadas durante el diseño del acelerador y con las predicciones de la teoría lineal de estabilidad longitudinal.

En primer lugar, se verifica la evolución energética de la partícula síncrona. Para una ganancia constante por brecha, la energía cinética después de $N_c$ etapas está determinada por

\begin{equation}
    T_{N_c}
    =
    T_0+
    N_cqV_{\mathrm{gap}}\cos(\phi_s).
\end{equation}

La comparación de esta expresión con los resultados del algoritmo de diseño permite comprobar la correcta implementación de la actualización energética.

En segundo lugar, se verifica la condición de sincronismo mediante la relación

\begin{equation}
    L_n=\beta_{s,n+1}\lambda.
\end{equation}

Esta condición garantiza que la partícula síncrona complete un período de radiofrecuencia durante el tránsito entre brechas consecutivas, conservando su fase de referencia.

Adicionalmente, se compara el avance de fase de las oscilaciones de sincrotrón obtenido mediante el seguimiento numérico con la predicción de la aproximación lineal, dada por la ecuación~\eqref{eq:mu}. Los resultados presentan una concordancia entre ambas estimaciones, dentro de la resolución del procedimiento discreto utilizado para determinar el avance de fase.

Desde el punto de vista cualitativo, las simulaciones también reproducen la dependencia de la estabilidad longitudinal respecto a la fase síncrona. Para protones, las fases situadas en el intervalo $-90^\circ<\phi_s<0^\circ$ presentan un comportamiento estable en la aproximación lineal, mientras que la operación en la cresta carece de una fuerza restauradora longitudinal de primer orden.

Finalmente, la utilización de una semilla aleatoria fija permite reproducir las mismas condiciones iniciales en ejecuciones sucesivas del programa. Sin embargo, la reproducibilidad de una realización particular no sustituye el análisis de la variabilidad estadística entre diferentes realizaciones del paquete.

\section{Limitaciones del modelo}

El modelo computacional permite estudiar la dinámica longitudinal del acelerador mediante un conjunto reducido de ecuaciones, pero incorpora simplificaciones que deben considerarse al interpretar los resultados.

En primer lugar, las brechas de aceleración se representan mediante una ganancia energética instantánea. Por tanto, no se considera la extensión espacial de las regiones aceleradoras ni la variación del campo eléctrico durante el tránsito dentro de ellas. En consecuencia, tampoco se incorpora un factor de tiempo de tránsito asociado a la distribución espacial del campo.

En segundo lugar, se considera una amplitud de tensión constante para todas las brechas. No se incluyen variaciones de amplitud entre cavidades ni fluctuaciones temporales de la señal de radiofrecuencia.

Asimismo, el modelo se limita al movimiento longitudinal de las partículas. No se consideran los grados de libertad transversales, la dinámica de focalización del haz ni los efectos de carga espacial. En un acelerador real, la estabilidad longitudinal puede estar acoplada a la dinámica transversal, por lo que se requieren elementos adicionales de focalización para controlar la divergencia del haz.

Por otra parte, el seguimiento se realiza mediante una discretización por celdas, utilizando la velocidad calculada después de cada brecha para aproximar el tiempo de tránsito por la celda correspondiente. Esta aproximación simplifica la integración de las ecuaciones de movimiento, pero no reproduce todos los detalles de la dinámica dentro de las cavidades.

Finalmente, la simulación del paquete considera una distribución estadística inicial de fase y energía, pero no incorpora una descripción completa de la distribución espacial del haz ni de sus interacciones colectivas.

Estas simplificaciones delimitan el alcance del modelo y deben tenerse presentes al extrapolar los resultados numéricos al comportamiento de un acelerador físico.

\section{Discusión}

Los resultados obtenidos permiten analizar la relación entre la condición de sincronismo, la estabilidad longitudinal y los parámetros de diseño de un acelerador lineal de protones con estructura tipo Alvarez.

La selección de la fase síncrona constituye uno de los aspectos fundamentales del diseño. Para $\phi_s=-30^\circ$, el acelerador alcanza una ganancia energética de aproximadamente $0.866\,\mathrm{MeV}$ por brecha y requiere 55 celdas para alcanzar la energía objetivo. Esta configuración proporciona una región de estabilidad longitudinal que permite mantener agrupada una fracción elevada del paquete de partículas.

Al modificar la fase síncrona se observa un compromiso entre la ganancia energética por brecha y las condiciones de estabilidad longitudinal. Las fases más alejadas de la cresta presentan una menor ganancia energética y, por tanto, requieren un mayor número de celdas. Sin embargo, también modifican la región longitudinal estable y la respuesta del paquete frente a las desviaciones iniciales de fase y energía.

La comparación entre las distintas fases síncronas muestra que la configuración de $-60^\circ$ proporciona la mayor fracción de partículas agrupadas entre los casos estudiados, aunque requiere una estructura considerablemente más larga que la configuración de $-30^\circ$. Por tanto, la elección de la fase síncrona no depende exclusivamente de maximizar la estabilidad, sino de establecer un compromiso entre la eficiencia energética y el comportamiento longitudinal del haz.

El análisis de las oscilaciones de sincrotrón permite identificar el mecanismo mediante el cual las desviaciones de fase y energía se acoplan durante el proceso de aceleración. La concordancia entre el avance de fase obtenido numéricamente y la predicción de la teoría lineal respalda la consistencia del modelo para pequeñas perturbaciones. Asimismo, la reducción de la amplitud de las oscilaciones es compatible con el amortiguamiento adiabático asociado al incremento de la energía relativista de las partículas.

Por otra parte, las simulaciones con errores de frecuencia muestran que el sincronismo es sensible a las variaciones de la señal de radiofrecuencia. A medida que aumenta el error relativo, se modifican las fases de llegada a las brechas y, en consecuencia, las ganancias energéticas experimentadas por las partículas. Para los parámetros considerados, los errores de hasta el $2\,\%$ conservan una fracción elevada del paquete dentro del criterio de agrupamiento establecido, mientras que un error del $4\,\%$ produce una reducción considerable de dicha fracción.

Finalmente, la comparación con electrones permite identificar diferencias importantes asociadas al régimen relativista. Debido a su menor masa en reposo, los electrones alcanzan velocidades próximas a la de la luz desde energías cinéticas relativamente bajas. En consecuencia, la deriva de fase por celda se reduce considerablemente cuando se utiliza una estructura de longitud constante aproximada por $\lambda$.

No obstante, es importante distinguir este comportamiento de la estabilidad longitudinal propiamente dicha. La reducción de la deriva de fase está relacionada con la proximidad de la velocidad de las partículas a $c$, mientras que la existencia de una fuerza restauradora longitudinal depende de la fase síncrona. En particular, para una operación en la cresta de radiofrecuencia, dicha fuerza desaparece en primer orden, independientemente de que la partícula acelerada sea un protón o un electrón.

\section{Conclusiones}

A partir del modelo computacional desarrollado y de los resultados obtenidos, se establecen las siguientes conclusiones:

\begin{enumerate}

    \item Se implementó un modelo discreto de aceleración por radiofrecuencia para una estructura de tubos de deriva tipo Alvarez. Mediante la condición de sincronismo y las relaciones relativistas se determinaron las longitudes de las celdas necesarias para acelerar protones desde $3\,\mathrm{MeV}$ hasta una energía objetivo de $50\,\mathrm{MeV}$. Para una fase síncrona de $-30^\circ$, se obtuvo una estructura de 55 celdas y una longitud total de $10.59\,\mathrm{m}$.

    \item El análisis del paquete de 3000 protones mostró que la fase síncrona influye significativamente en la estabilidad longitudinal. Para $\phi_s=-30^\circ$, el $99.3\,\%$ de las partículas satisface el criterio de agrupamiento establecido, mientras que para la operación en la cresta de radiofrecuencia esta fracción disminuye considerablemente. Estos resultados evidencian la importancia de seleccionar adecuadamente la fase síncrona para mantener la concentración longitudinal del haz.

    \item La simulación de las oscilaciones de sincrotrón presentó concordancia con las predicciones de la teoría lineal para pequeñas perturbaciones. Se observó una disminución del avance de fase por celda y de la amplitud de las oscilaciones conforme aumenta la energía relativista de la partícula, comportamiento consistente con el amortiguamiento adiabático.

    \item El análisis de errores en la frecuencia de radiofrecuencia mostró que las desviaciones respecto a la frecuencia nominal modifican la evolución de fase y la energía final del paquete. Para los parámetros estudiados, se conservó una elevada fracción de partículas agrupadas ante errores de frecuencia de hasta el $2\,\%$, mientras que un error del $4\,\%$ produjo una reducción considerable de la captura longitudinal.

    \item La comparación con electrones permitió identificar la influencia del régimen ultrarrelativista sobre la evolución de fase. La proximidad de su velocidad a la velocidad de la luz reduce la deriva de fase durante el tránsito por celdas de longitud constante. Sin embargo, la ausencia de una fuerza restauradora longitudinal de primer orden en la operación de cresta está determinada por la fase síncrona y no exclusivamente por la naturaleza de la partícula.

\end{enumerate}
